% Part: propositional-logic
% Chapter: syntax-and-semantics
% Section: formulas

\documentclass[../../../include/open-logic-section]{subfiles}

\begin{document}

\olfileid{pl}{syn}{fml}

\olsection{قضیاتی \usetoken{P}{formula}}

قضیاتی منطق دے !!^{formula}s
\emph{!!{propositional
    variable}s} توں
\iftag{prvFalse}{\iftag{prvTrue}{،}{ تے} قضیاتی مستقل~$\lfalse$}{}
\iftag{prvTrue}{\iftag{prvFalse}{ تے}{} قضیاتی مستقل~$\ltrue$}{}
\emph{منطقی رابطیاں} نال بنائے جاندے نیں۔

\begin{enumerate}
\item !!^a{denumerable} سیٹ~$\PVar$، جیہڑا !!{propositional variable}s
  $\Obj p_0$، $\Obj p_1$، \dots دا اے۔
\tagitem{prvFalse}{!!{falsity} دا قضیاتی مستقل~$\lfalse$۔}{}
\tagitem{prvTrue}{!!{truth} دا قضیاتی مستقل~$\ltrue$۔}{}
\item منطقی رابطے:
  \startycommalist
  \iftag{prvNot}{\ycomma $\lnot$ (نفی)}{}%
  \iftag{prvAnd}{\ycomma $\land$ (عطف)}{}%
  \iftag{prvOr}{\ycomma $\lor$ (فصل)}{}%
  \iftag{prvIf}{\ycomma $\lif$ (شرطیہ، !!{conditional})}{}%
  \iftag{prvIff}{\ycomma $\liff$ (دو شرطیہ، !!{biconditional})}{}%
\item رموزِ اوقاف: (، )، تے کوما۔
\end{enumerate}

اسیں قضیاتی منطق دی ایس زبان نوں $\Lang L_0$ نال ظاہر کردے آں۔

\iftag{defNot,defOr,defAnd,defIf,defIff,defTrue,defFalse,defEx,defAll}{%
اُتے متعارف کرائے ابتدائی رابطیاں دے نال نال، اسیں ایہ
\emph{تعریف شدہ} علامات وی ورتدے آں:
\startycommalist
  \iftag{defNot}{\ycomma $\lnot$ (نفی)}{}%
  \iftag{defAnd}{\ycomma $\land$ (عطف)}{}%
  \iftag{defOr}{\ycomma $\lor$ (فصل)}{}%
  \iftag{defIf}{\ycomma $\lif$ (شرطیہ، !!{conditional})}{}%
  \iftag{defIff}{\ycomma $\liff$ (دو شرطیہ، !!{biconditional})}{}%
  \iftag{defFalse}{\ycomma $\lfalse$ (کذب، !!{falsity})}{}%
% PNBPL-001 ماخذ دی درستی: اندرلے defTrue شرطی حکم دی خالی تیسری شاخ ماخذ وچ رہ گئی اے؛ ایتھے اوہ عین خالی شاخ شامل اے۔
  \iftag{defTrue}{\ycomma $\ltrue$ (صدق، !!{truth})}{}}{}۔

\begin{tagblock}{defNot,defOr,defAnd,defIf,defIff,defTrue,defFalse,defEx,defAll}
\begin{explain}
تعریف شدہ علامت باقاعدہ طور تے زبان دا حصہ نہیں ہُندی، بلکہ اک غیر رسمی
مخفف دے طور تے متعارف کرائی جاندی اے: ایہ سانوں اوہناں فارمولیاں نوں
مختصر لکھن دی سہولت دیندی اے جیہڑے صرف ابتدائی علامات نال لکھے جان تاں
خاصے لمّے ہو جان۔ ایہ صاف طور تے اک فائدہ اے۔ پر ایہ توں وڈا فائدہ ایہ
اے کہ ثبوت چھوٹے ہو جاندے نیں۔ جے کوئی علامت ابتدائی ہووے، تاں ثبوتاں
وچ اوہنوں وکھرے طور تے ورتنا پیندا اے۔ ایس لئی جِنّیاں زیادہ ابتدائی
علامتاں ہون، ساڈے ثبوت اونّے ای لمّے ہون گے۔
\end{explain}
\end{tagblock}

% Alternate symbols

\begin{intro}
ہو سکدا اے تُسی اُتّے ساڈی ورتی اصطلاح تے علامات توں وکھری اصطلاح تے
علامتاں نال واقف ہووو۔ منطق دیاں کتاباں (تے استاد) ``نفی'' لئی عام طور
تے $\sim$، $\neg$ تے~!\ ورتدے نیں، تے ``عطف'' لئی $\wedge$،
$\cdot$ تے $\&$۔ ``شرطیہ'' یا ``لزوم'' لئی عام ورتیاں جان والیاں
علامتاں $\rightarrow$، $\Rightarrow$ تے $\supset$ نیں۔
\iftag{prvIff,defIff}{``دو شرطیہ''، ``دو طرفہ لزوم'' یا ``(مادی) ہم
  ارزی'' دیاں علامات $\leftrightarrow$، $\Leftrightarrow$ تے
  $\equiv$ نیں۔}{}
\iftag{prvFalse,defFalse}{$\lfalse$ دی علامت نوں وکھ وکھ تھاواں اُتے
  ``کذب''، ``کاذب''، ``محال'' یا ``زیریں'' آکھیا جاندا اے۔}{}
\iftag{prvTrue,defTrue}{$\ltrue$ دی علامت نوں وکھ وکھ تھاواں اُتے
  ``صدق''، ``صادق'' یا ``بالائی'' آکھیا جاندا اے۔}{}
\end{intro}

\begin{defn}[فارمولا]
\ollabel{defn:formulas}
قضیاتی منطق دے \emph{!!{formula}s} دا سیٹ~$\Frm[L_0]$ استقرائی طور تے
اینج تعریف کیتا جاندا اے:
\begin{enumerate}
\tagitem{prvFalse}{$\lfalse$ اک ایٹمی !!{formula} اے۔}{}

\tagitem{prvTrue}{$\ltrue$ اک ایٹمی !!{formula} اے۔}{}

\item ہر !!{propositional variable}~$\Obj p_i$ اک ایٹمی
  !!{formula} اے۔

\tagitem{prvNot}{جے $!A$، !!a{formula} اے، تاں $\lnot !A$ وی
  !!a{formula} اے۔}{}

\tagitem{prvAnd}{جے $!A$ تے $!B$، !!{formula}s نیں، تاں $(!A \land
  !B)$ وی !!a{formula} اے۔}{}

\tagitem{prvOr}{جے $!A$ تے $!B$، !!{formula}s نیں، تاں $(!A \lor !B)$
  وی !!a{formula} اے۔}{}

\tagitem{prvIf}{جے $!A$ تے $!B$، !!{formula}s نیں، تاں $(!A \lif !B)$
  وی !!a{formula} اے۔}{}

\tagitem{prvIff}{جے $!A$ تے $!B$، !!{formula}s نیں، تاں $(!A \liff !B)$
  وی !!a{formula} اے۔}{}

\tagitem{limitClause}{ہور کوئی شے !!a{formula} نہیں اے۔}{}
\end{enumerate}
\end{defn}

\begin{explain}
!!{formula}s دی تعریف اک
\emph{استقرائی تعریف} اے۔ بنیادی طور تے اسیں !!{formula}s دا سیٹ
لامتناہی مرحلیاں وچ بناؤندے آں۔ پہلے مرحلے وچ اسیں سارے ایٹمی
فارمولیاں نوں فارمولے قرار دیندے آں؛ ایہ تعریف دیاں پہلیاں کجھ صورتاں،
یعنی
\iftag{prvTrue}{$\ltrue$، }{}%
\iftag{prvFalse}{$\lfalse$، }{}%
$\Obj p_i$ والیاں صورتاں دے مطابق اے۔ ایس طرح ``ایٹمی !!{formula}''
دا مطلب ایس روپ دا کوئی وی !!{formula} اے۔

تعریف دیاں ہور صورتاں پہلے بنائے !!{formula}s توں نویں !!{formula}s
بناؤن دے قاعدے دیندیاں نیں۔ دوجے مرحلے وچ اسیں ایہناں قاعدیاں نال ایٹمی
!!{formula}s توں !!{formula}s بنا سکدے آں۔ تیجے مرحلے وچ اسیں ایٹمی
فارمولیاں تے دوجے مرحلے وچ ملے فارمولیاں توں نویں فارمولے بناؤندے آں،
تے ایویں ای اگے۔ !!{formula} اوہو شے اے جیہڑی آخرکار کسے ایہو جہے
مرحلے اُتے بنائی جاوے، تے ہور کوئی شے نہیں۔
\end{explain}

فارمولا $(!B \ast !C)$، جیہڑا $!B$، $!C$ توں دو-تھاں
رابطے~$\ast$ نال بنیا اے، لکھدے ہویاں، اسیں اکثر سب توں باہری قوسین
دی جوڑی چھڈ کے صرف
$!B \ast !C$ لکھاں گے۔

\begin{tagblock}{defNot,defOr,defAnd,defIf,defIff,defTrue,defFalse,defEx,defAll}
\begin{defn}
تعریف شدہ عامل ورت کے بنائے فارمولیاں نوں اینج سمجھیا جاوے:

\begin{tagenumerate}{defTrue,defFalse,defNot,defOr,defAnd,defIf,defIff,defEx,defAll}

\tagitem{defTrue}{$\ltrue$،
  \iftag{prvFalse}{$\lnot\lfalse$}{$(!A \lor \lnot !A)$، کسے مقرر
    ایٹمی !!{formula}~$!A$ لئی} دا مخفف اے۔}{}

\tagitem{defFalse}{$\lfalse$،
  \iftag{prvTrue}{$\lnot\ltrue$}{$(!A \land \lnot !A)$، کسے مقرر
    ایٹمی !!{formula}~$!A$ لئی} دا مخفف اے۔}{}

\tagitem{defNot}{$\lnot !A$، $!A \lif \lfalse$ دا مخفف اے۔}{}

\tagitem{defOr}{$!A \lor !B$،
  \iftag{prvAnd}{$\lnot(\lnot !A \land \lnot !B)$}{$\lnot !A \lif
    !B$} دا مخفف اے۔}{}

\tagitem{defAnd}{$!A \land !B$،
  \iftag{prvOr}{$\lnot(\lnot !A \lor \lnot !B)$}{$\lnot (!A \lif
    \lnot !B)$} دا مخفف اے۔}{}

% PNBPL-002 ماخذ دی درستی: پہلی متبادل عبارت دے آخری قوس دی شروع والی قوس ماخذ وچ رہ گئی اے؛ ایتھے پوری جوڑی شامل اے۔
\tagitem{defIf}{$!A \lif !B$،
  \iftag{prvOr}{$(\lnot !A \lor !B)$}{$\lnot (!A \land \lnot !B)$}
  دا مخفف اے۔}{}

\tagitem{defIff}{$!A \liff !B$، $(!A \lif !B) \land (!B
  \lif !A)$ دا مخفف اے۔}{}
\end{tagenumerate}
\end{defn}
\end{tagblock}

\begin{defn}[نحوی یکسانیت]
علامت $\ident$، علامتاں دیاں لڑیاں وچ نحوی یکسانیت ظاہر کردی اے، یعنی
$!A \ident !B$ اُدوں تے صرف اُدوں جدوں $!A$ تے $!B$ علامتاں دیاں
اکّو لمبائی دیاں لڑیاں ہون تے ہر مقام اُتے اکّو علامت رکھدیاں ہون۔
\end{defn}

علامت $\ident$ دے دونیں پاسے جوڑ نال ملیاں لڑیاں آ سکدیاں نیں؛ مثلاً
$!A \ident (!B \lor !C)$ دا مطلب اے: علامتاں دی لڑی~$!A$ اوہی لڑی
اے جیہڑی، ایس ترتیب نال، اک کھلدی قوس، لڑی $!B$، علامت $\lor$،
لڑی~$!C$ تے اک بند ہُندی قوس نوں جوڑ کے ملدی اے۔ جے ایہ گل اے، تاں
سانوں پتا اے کہ $!A$ دی پہلی علامت اک کھلدی قوس اے، $!A$ وچ $!B$ اک
ذیلی لڑی دے طور تے شامل اے (دوجی علامت توں شروع ہوندی اے)، اوس ذیلی
لڑی دے بعد $\lor$ آؤندا اے، وغیرہ۔

\end{document}
